\begin{abstract}
Percussive one-shots range from pitched drums to cymbals with dense modal content, and sound designers often want to adjust the timbre of an existing sample rather than generate a new one.
Differentiable analysis--synthesis models such as DrumBlender reconstruct a given sample through interpretable components, but their low-level parameters offer no high-level control, and their constant-Q analysis struggles with dense high-frequency modes.
This paper builds on DrumBlender and reconstructs a one-shot as the sum of jointly trained modal, noise, and transient components.
Sinusoidal tracks from a hybrid constant-Q and short-time Fourier analysis receive bounded corrections, noise bands following NoiseBandNet are conditioned on the corrected modal parameters, and a small temporal convolutional network adds a transient computed from the modal waveform.
After training, a sparse autoencoder fitted to the shared representation provides label-free controls whose edits start from the reconstruction of the given sample and act through the synthesis parameters.
\textbf{[TBD: one or two sentences on the measured reconstruction and control results.]}
\end{abstract}

\begin{keywords}
differentiable digital signal processing; percussion synthesis; timbre control; modal analysis; sparse autoencoders
\end{keywords}
